\documentclass[12pt]{article}
\usepackage[utf8]{inputenc}
\usepackage{amsmath, amssymb}
\usepackage{listings}
\usepackage{xcolor}
\usepackage{enumitem}
\usepackage[margin=1in]{geometry}
\usepackage{hyperref}
\usepackage[skins]{tcolorbox}


\usepackage{changepage}  % For margin adjustments
\newenvironment{codemargins}{\begin{adjustwidth}{-2em}{-2em}}{\end{adjustwidth}}


\newcommand{\addmedskip}{\addvspace{\medskipamount}}
\newcommand{\addbigskip}{\addvspace{\bigskipamount}}


\definecolor{codegreen}{rgb}{0,0.6,0}
\definecolor{codegray}{rgb}{0.5,0.5,0.5}
\definecolor{codepurple}{rgb}{0.58,0,0.82}
\definecolor{backcolour}{rgb}{0.95,0.95,0.92}

\lstdefinestyle{mystyle}{
    backgroundcolor=\color{backcolour},   
    commentstyle=\color{codegreen},
    keywordstyle=\color{magenta},
    numberstyle=\tiny\color{codegray},
    stringstyle=\color{codepurple},
    basicstyle=\footnotesize\ttfamily,
    breakatwhitespace=false,         
    breaklines=true,                 
    breakautoindent=true,            % <-- enable auto-indent for wrapped lines
    breakindent=2em,                 % <-- indentation for wrapped lines
    postbreak=\mbox{\textcolor{red}{$\hookrightarrow$}\space}, % marker for wrap
    captionpos=b,                    
    keepspaces=true,                 
    numbers=left,                    
    numbersep=5pt,                  
    showspaces=false,                
    showstringspaces=false,
    showtabs=false,                  
    tabsize=2,
    escapeinside={(*@}{@*)}
}

\lstdefinelanguage{Solidity}{
    keywords={contract, function, if, else, return, returns, for, while, true, false, public, private, internal, external, payable, view, pure, virtual, override, import, from, as, struct, enum, mapping, address, uint, bytes, bool, string, event, emit, modifier, require, assert, revert, new, delete, try, catch, break, continue, constant, immutable, anonymous, indexed, storage, memory, calldata, constructor, fallback, receive},
    keywordstyle=\color{blue}\bfseries,
    ndkeywords={pragma, solidity, is, using, library},
    ndkeywordstyle=\color{darkgray}\bfseries,
    identifierstyle=\color{black},
    sensitive=false,
    comment=[l]{//},
    morecomment=[s]{/*}{*/},
    commentstyle=\color{purple}\ttfamily,
    stringstyle=\color{red}\ttfamily,
    morestring=[b]',
    morestring=[b]"
}

\lstset{style=mystyle}


%%%%%%%%% PROBLEM ENVIRONMENT %%%%%%%%%%%%
\newcounter{problemgroup}
\setcounter{problemgroup}{1} % Start at 1
\newcounter{problemnum}[problemgroup]
\newenvironment{problem}{%
    \addbigskip 
    \setcounter{partnum}{0} % Keeping your part reset
    \noindent
    \stepcounter{problemnum}%
    \textbf{Problem \arabic{problemgroup}.\arabic{problemnum}.\ }%
}{%
    \par\addbigskip
}
\newcommand{\nextproblemsection}{%
    \stepcounter{problemgroup}%
}

\newcounter{partnum}
\setcounter{partnum}{0}
\newenvironment{ppart}
     {\addmedskip
      \noindent\stepcounter{partnum}\textbf{(\alph{partnum})}\ }
     {\par\addbigskip}

%%%%%% SOLUTION ENVIRONMENT %%%%%%%%
\newtcolorbox{solution}{
    enhanced,
    colback=white,
    colframe=black,
    arc=0mm,
    boxrule=0.5pt,
    left=10pt, right=10pt, top=10pt, bottom=10pt,
    before=\vspace{1em},
    after=\vspace{1em},
    height=2.5in,
}

\newtcolorbox{bigsolution}{
    enhanced,
    colback=white,
    colframe=black,
    arc=0mm,
    boxrule=0.5pt,
    left=10pt, right=10pt, top=10pt, bottom=10pt,
    before=\vspace{1em},
    after=\vspace{1em},
    height=4.3in,
}

\newtcolorbox{verybigsolution}{
    enhanced,
    colback=white,
    colframe=black,
    arc=0mm,
    boxrule=0.5pt,
    left=10pt, right=10pt, top=10pt, bottom=10pt,
    before=\vspace{1em},
    after=\vspace{1em},
    height=7.3in,
}

\newif\ifshowsolutions
\IfFileExists{26s/hw5/hw5_solsNOTGIVEN.tex}{
    \showsolutionstrue
    \input{26s/hw5/hw5_sols.tex}
}{
    \showsolutionsfalse
}
\ifshowsolutions
    % Do nothing, the commands are loaded from the file
\else
    % Define dummy commands so the TeX still compiles for students
    \newcommand{\HwFiveSolOneOne}{}
    \newcommand{\HwFiveSolOneTwo}{}
    \newcommand{\HwFiveSolOneThree}{}
    \newcommand{\HwFiveSolOneFour}{}
    \newcommand{\HwFiveSolOneFive}{}
    \newcommand{\HwFiveSolOneSix}{}
    \newcommand{\HwFiveSolOneSeven}{}
    \newcommand{\HwFiveSolOneEight}{}
    \newcommand{\HwFiveSolTwoOne}{}
    \newcommand{\HwFiveSolTwoTwo}{}
    \newcommand{\HwFiveSolTwoThree}{}
    \newcommand{\HwFiveSolTwoFour}{}
    \newcommand{\HwFiveSolTwoFive}{}
    \newcommand{\HwFiveSolTwoSix}{}
    \newcommand{\HwFiveSolTwoSeven}{}
    \newcommand{\HwFiveSolThreeOne}{}
    \newcommand{\HwFiveSolThreeTwo}{}
    \newcommand{\HwFiveSolThreeThree}{}
    \newcommand{\HwFiveSolThreeFour}{}
\fi
%%%%%%%%% Done Handling Solutions %%%%%%%%%%%%%



\begin{document}

\begin{center}
\Large{\textbf{CS595 Homework 6}}\\
\bigskip
\end{center}

\section*{Section 1}


\begin{problem} \textbf{(5 points).}
A student named Stu is implementing the mixer from this assignment. He notices that
anyone can see the commitments of all deposits by looking at the blockchain. Since all the
commitments are public, anyone can compute a Merkle path from each commitment to the
root. Stu reasons that because all paths can be computed by anyone, he can make the path
in the withdraw.nr zero-knowledge proof be a public input. What is the problem with
Stu’s scheme?
\end{problem}

\begin{solution}
    \HwFiveSolOneEight The whole point of this system is to preserve the anonymity of transactions. If the hashpath to deposit is public, then people will be able to trace withdrawals to deposits.
    However, if the index is still not public an attacker would only be able to determine whether a certain commitment or its sibling is the withdrawn commitment which would still be damaging to anonymity.
    Although an attacker can compute any hash path, they still wouldn't know which specific hash path was used, which preserves privacy.
\end{solution}


\nextproblemsection
\clearpage

\section*{Section 2}

\begin{problem}
Stu has another idea. Instead of making the Merkle path in the withdrawal zero-
knowledge proof public, he just removes it from the zero-knowledge proof entirely. He writes
a withdraw.nr function that takes in r and c as a private input, id as a public input, and
checks that PedersenHash(id, r) = c. What is the problem with this scheme?
\end{problem}

\begin{solution}
    \HwFiveSolOneEight
\end{solution}

\section*{Acknowledgment of AI and external resources}

\begin{problem} (0 points)
    Please provide links to external references you used for this assignment and any prompts, session links or session copies for substantial use of generative AI, or attest that you did not use any external resources.
\end{problem}

\begin{bigsolution}
    I did not use generative AI in any substantial way for this problem set.
    \vspace{2in}
\end{bigsolution}


\end{document}